\documentclass[11pt]{article}

\usepackage[margin=1in]{geometry}
\usepackage{amsmath,amssymb,amsthm}
\usepackage{booktabs}
\usepackage{microtype}

\newtheorem{theorem}{Theorem}
\newtheorem{remark}{Remark}

\newcommand{\oracle}{\widehat{O}}
\newcommand{\eps}{\varepsilon}

\title{The mass spectrum of the IIB flux landscape\\ carries no geometric fingerprint}
\author{Rob Mitchell}
\date{}

\begin{document}
\maketitle

\begin{abstract}
A central hope of the ``structure mining'' program is that some low-energy
observable of a flux vacuum carries a \emph{fingerprint}---information about the
flux that generated it, beyond its overall scale---so that the compactification
could be pinned down from observables. We test the last candidate: the mass
spectrum. On $133$ supersymmetric IIB flux vacua (JAXVacua, $h^{1,2}=2$, natural
ISD measure) we compute the degree-$0$ spectrum---the six eigenvalues of the
complex-structure moduli and the axio-dilaton---and measure how predictable each
eigenvalue is from the flux. The raw eigenvalues are predictable
(ridge-on-quadratic CV $r = 0.50$--$0.62$, heaviest slot significant at
$z=5.5$), but the \emph{scale-invariant} ratios $m^2_i / m^2_{\rm max}$ are not
($r < 0.25$). The conclusion: the mass spectrum's \emph{shape} is statistically
uniform in flux---the fingerprint is not there---while its \emph{overall scale}
is flux-determined through a quadratic form distinct from the D3 tadpole. This
closes, in the flux coordinates, the ``masses carry the fingerprint'' hypothesis,
and leaves the constraint (the tadpole) as the only learnable structure found to
date.
\end{abstract}

\section{Introduction}

The string landscape is $\sim 10^{500}$ vacua, and no experiment distinguishes
most of them. The problem of \emph{verification} (does some vacuum match our
universe?) is the subject of a companion paper \cite{pacland}; here we test one
specific hypothesis that has guided the search for structure in the landscape:
that \emph{the mass spectrum is the fingerprint}---the low-energy observable
that carries information about the flux, so that measuring the masses would pin
the compactification.

This hope is natural. The superpotential $W_0$, the string coupling $g_s$, and
the gauge couplings have all been measured to be statistically uniform in the
flux coordinates (cross-validated $r < 0.25$) \cite{pacland}: they are random
draws that carry no information about the flux quanta that produced them. The
mass spectrum was the one observable not yet measured, and it was the natural
candidate for the fingerprint, because it depends on the full geometry through
the stabilized moduli rather than on a single closed form.

We compute it, and we find the hope is not borne out.

\section{Method}

\paragraph{The sample.} We use JAXVacua \cite{jaxvacua} on a bundled $h^{1,2}=2$
flux EFT (``KS'' model). At ${\rm maximum\_degree}=0$ (no instanton corrections)
the K\"ahler moduli are flat directions, so the mass matrix has $6$ physical
eigenvalues: $4$ from the two complex-structure moduli and $2$ from the
axio-dilaton. Critical points are sampled with the natural ISD measure
(${\tt isd\_mode}={\tt ISD-}$, Gaussian flux prior); supersymmetric vacua are
selected by $D W = 0$ and $\operatorname{Im}\tau > 0$. We obtain $133$ SUSY
vacua from $964$ critical points.

\paragraph{The test.} For each eigenvalue $m^2_i$ we measure the cross-validated
ridge regression of $m^2_i$ on the degree-$2$ monomial basis of the $12$ flux
coordinates ($91$ features), reporting the Pearson $r$ between held-out
predictions and truth. This is the same joint learnability test used throughout
\cite{pacland}: $r$ near $1$ means deterministic structure (the tadpole control),
$r$ near $0$ means statistically uniform.

\paragraph{Scale vs.\ shape.} A nonzero $r$ can arise in two ways: the
\emph{scale} of the spectrum (the overall curvature, set by the potential) or the
\emph{shape} (the ratios $m^2_i / m^2_j$). To separate them we re-run the test on
the scale-invariant ratios $m^2_i / m^2_{\rm max}$. If the shape carries the
fingerprint, the ratios are learnable; if only the scale does, the raw values are
learnable but the ratios are not.

\section{Results}

\begin{center}
\begin{tabular}{lcccccc}
\toprule
 & \multicolumn{6}{c}{eigenvalue slot (sorted ascending)} \\
 & 0 & 1 & 2 & 3 & 4 & 5 \\
\midrule
raw $m^2_i$ & $0.03$ & $0.37$ & $0.50$ & $0.53$ & $0.58$ & $\mathbf{0.62}$ \\
ratio $m^2_i / m^2_5$ & $0.09$ & $0.16$ & $0.21$ & $0.15$ & $0.08$ & --- \\
\bottomrule
\end{tabular}
\end{center}

The raw eigenvalues are learnable: the three heaviest slots cross the
$r = 0.5$ threshold, and the heaviest ($m^2_5$) is significant at $z = 5.5$
against a permutation null (shuffled flux labels; null mean $\approx 0.001$,
std $\approx 0.11$). The scale-invariant ratios are \emph{not} learnable: all
$r < 0.25$, inside the null. The heaviest eigenvalue correlates with the D3
tadpole at only $r = 0.12$, so the learned scale is a quadratic form of the flux
\emph{distinct} from the tadpole, not the tadpole itself.

\section{Conclusion}

\begin{remark}
\textbf{The mass spectrum's shape is statistically uniform in flux; its scale is
flux-determined.} In the IIB flux corner, on $133$ SUSY vacua, the degree-$0$
mass eigenvalues carry information about the flux only through their overall
scale. The scale-invariant ratios---the part that would constitute a geometric
fingerprint---are random ($r < 0.25$). The fingerprint is not in the mass
spectrum.
\end{remark}

This closes the ``masses carry the fingerprint'' hypothesis in the flux
coordinates, and completes the pattern: $W_0$, $g_s$, the gauge couplings, and now
the mass-spectrum shape are all statistically uniform in flux. The only learnable
structure found to date is the \emph{constraint}---the D3 tadpole, which is the
exact symplectic pairing. The constraint-first program of \cite{pacland} therefore
survives its first decisive falsification attempt: the landscape's navigable
structure is what it \emph{forbids}, not what it \emph{selects}.

\paragraph{Scope, named.} (i) This is the degree-$0$ spectrum; the
instanton-corrected K\"ahler-moduli masses remain computationally infeasible and
could in principle differ. (ii) ``Random'' means \emph{in the flux coordinates};
the masses could still carry structure in the moduli coordinates, which we have
not tested. (iii) It is a measurement on one model ($h^{1,2}=2$) at
$133$ vacua, not a theorem.

\begin{thebibliography}{9}
\bibitem{pacland} R.~Mitchell, ``Probably Approximately Correct String Theory:
A Verification Discipline for the Landscape'' (companion paper).
\bibitem{jaxvacua} A.~Dubey, S.~Krippendorf, A.~Schachner, ``JAXVacua''
(arXiv:2306.06160).
\bibitem{denef} F.~Denef and M.~R.~Douglas, ``Distributions of flux vacua,''
JHEP \textbf{0405} (2004) 072.
\end{thebibliography}

\end{document}
